\chapter{Functions}
\label{ch:functions}

\begin{goals}
\begin{itemize}
  \item What a function is and why we divide a program into functions
  \item Defining a function and calling it
  \item Sending values into a function with parameters
  \item Getting a result back from a function with \code{ret}
  \item Default values for parameters, and the scope of variables
\end{itemize}
\end{goals}

\section{One name for one job}

Since the first chapter we have been working with one function:
\code{func main}. Now it is time to see what a ``function'' really is. A
\textbf{function} is a block of instructions that has a name. Wherever we
write that name (this is called \textbf{calling} the function), the
instructions inside the function run. \code{main} is the same thing, except
that Salam calls it by itself, at the start of the program.

Suppose a program has to print a separator line in several places. We could
write \code{println "--------------------"} each time, or we could define
the job once in a function and call its name wherever it is needed:

\program{Defining and calling a function}
\begin{salamcode}
func separator line():
    println "--------------------"
end

func main:
    separator line()
    println "Shopping list"
    separator line()
    println "bread, cheese, milk"
    separator line()
end
\end{salamcode}

\begin{salamout}
--------------------
Shopping list
--------------------
bread, cheese, milk
--------------------
\end{salamout}

A function definition starts with the word \kwd{func}, then the name of the
function, then a pair of parentheses and a colon, and finally \code{end}.
Like a variable name, a function name can be made of several words. A call
is the name of the function followed by parentheses:
\code{separator line()}. The parentheses are required; if you leave them
out, Salam reports an error: ``routine \code{'separator line'} is not
called: write \code{'separator line()'} to call it''.

Two things to notice. First, we defined the function \emph{outside}
\code{main}, at the top of the program; functions are always defined at the
top level of the program, side by side, never inside one another. Second,
the order of the definitions does not matter; we could have written
\code{separator line} below \code{main} just as well. Salam reads all the
definitions first and then starts from \code{main}.

\begin{note}[Why functions?]
There are three main reasons. \textbf{No repetition}: a job that is needed
in several places is written once, and if it ever has to change, it changes
in one place. \textbf{Naming}: reading \code{separator line()} is far
easier than reading a row of dashes in the middle of the code; the name
states the intent. \textbf{Dividing the work}: we break a large problem
into several small functions and write and test each one separately.
\end{note}

\section{Parameters: sending values into a function}

The function \code{separator line} always does the same fixed job. Most
functions, however, need some information to do their work: the name of the
person to greet, the number to square. We send this information into the
function through \textbf{parameters}. Parameters are written inside the
parentheses of the function definition, and for each one we must also state
its \emph{type}:

\program{A function with a parameter}
\begin{salamcode}
func greet(name: str):
    println "Hello", name + "! Welcome."
end

func main:
    greet("Sara")
    greet("Reza")
    friend := "Mina"
    greet(friend)
end
\end{salamcode}

\begin{salamout}
Hello Sara! Welcome.
Hello Reza! Welcome.
Hello Mina! Welcome.
\end{salamout}

\code{name: str} says that this function wants one parameter called
\code{name} of type string. In each call, the value we write inside the
parentheses (called the \textbf{argument}) is put into the box \code{name},
and the function works with it. An argument can be a value written directly,
a variable, or even a calculation.

A function can have several parameters, separated by commas. When it is
called, the arguments reach the parameters in the same order:

\program{Several parameters}
\begin{salamcode}
func mark row(count: int, mark: str):
    mut line := ""
    repeat count:
        line += mark
    end
    println line
end

func main:
    mark row(5, "*")
    mark row(3, "#")
    mark row(8, "=")
end
\end{salamcode}

\begin{salamout}
*****
###
========
\end{salamout}

\begin{warn}
The type of an argument must match the type of the parameter. The call
\code{mark row("5", "*")} is an error, because \code{"5"} is a string and
the first parameter wants an integer. The number of arguments must be right
as well: neither fewer nor more.
\end{warn}

\section{Returning: getting a result from a function}

The functions above did something (printing) but gave nothing back to their
caller. Many functions have to \emph{give} a result: an area, a final price,
the answer to a question. Two things are needed for this: we declare the
type of the result after the parentheses, and inside the function we hand
the result back with \kwd{ret}:

\program{A function with a return value}
\begin{salamcode}
func rectangle area(length: int, width: int): int:
    ret length * width
end

func main:
    println "Area of the room:", rectangle area(4, 5)
    println "Area of the yard:", rectangle area(12, 8)
    total := rectangle area(4, 5) + rectangle area(12, 8)
    println "Total:", total
end
\end{salamcode}

\begin{salamout}
Area of the room: 20
Area of the yard: 96
Total: 116
\end{salamout}

Look at the head of the function:
\code{func rectangle area(length: int, width: int): int:}. After the
parentheses come a colon and the type of the result (\code{int}), and then
the usual colon that opens the block. The instruction
\code{ret length * width} works out the value, hands it to the caller, and
the function finishes right there.

The most important point is that a call to such a function is itself a
\emph{value}, and it can be put anywhere a value can go: in a
\code{println}, in a calculation, in the declaration of a variable, or even
as the argument of another function.

\subsection{A function that says yes or no}

A function that gives a boolean result is very handy in conditions:

\program{A boolean function}
\begin{salamcode}
func is even(n: int): bool:
    ret (n % 2) == 0
end

func main:
    repeat 1 to 6 in n:
        if is even(n):
            println n, "even"
        else:
            println n, "odd"
        end
    end
end
\end{salamcode}

\begin{salamout}
1 odd
2 even
3 odd
4 even
5 odd
6 even
\end{salamout}

The line \code{if is even(n):} reads almost like English. This is one of
the beauties of naming functions well. It is a good custom to name boolean
functions as questions: \code{is even}, \code{is valid}, \code{is empty}.

\subsection{Several returns in one function}

A function may have \code{ret} in several places. As soon as the first
\code{ret} runs, the function is finished and the remaining lines do not
run. This feature shortens chains of conditions:

\program{Several returns}
\begin{salamcode}
func grade label(grade: int): str:
    if grade >= 18: ret "Excellent" end
    if grade >= 15: ret "Good" end
    if grade >= 10: ret "Acceptable" end
    ret "Fail"
end

func main:
    println grade label(19)
    println grade label(16)
    println grade label(12)
    println grade label(7)
end
\end{salamcode}

\begin{salamout}
Excellent
Good
Acceptable
Fail
\end{salamout}

With a grade of 16, the first condition is false, the second is true, and
the function returns with ``Good''; the lines after it are never seen.
Compare this with the same job in chapter \ref{ch:conditions}, which we
wrote with \code{else}.

\begin{warn}
Salam is as strict with functions as it is with variables: a function that
is defined but never called anywhere is an error. Also, a function that has
declared a return type must return something on \emph{every} path. If you
delete the final \code{ret "Fail"}, Salam gives a \emph{warning} that the
function can reach the end of its body without returning a value. The
program runs despite the warning, but for a grade of 7 the function no
longer returns a label, and \code{null} is printed instead; so take warnings
as seriously as errors. A function like \code{greet}, which only does
something and gives nothing back, does not need \code{ret}.
\end{warn}

\subsection{A function with floating-point numbers}

\program{Temperature conversion as a function}
\begin{salamcode}
func fahrenheit(celsius: f64): f64:
    ret ((celsius * 9.0) / 5.0) + 32.0
end

func main:
    println "Water freezes:", fahrenheit(0.0)
    println "Body temperature:", fahrenheit(37.0)
    println "Water boils:", fahrenheit(100.0)
end
\end{salamcode}

\begin{salamout}
Water freezes: 32
Body temperature: 98.6
Water boils: 212
\end{salamout}

The formula we wrote inside a loop in chapter \ref{ch:loops} now has a name
and can be used wherever it is needed. Notice that we wrote the arguments
with a decimal point (\code{0.0}), because the parameter is floating-point.
Salam would accept \code{0} too and convert it by itself, but \code{0.0}
shows the intent more clearly.

\section{A default value for a parameter}

Sometimes a parameter has a particular value most of the time and changes
only now and then. We can give it a \textbf{default value} in the function
definition; the caller may then leave that argument out:

\program{A parameter with a default value}
\begin{salamcode}
func greet(name: str, prefix: str = "Hello"):
    println prefix, name
end

func main:
    greet("Sara")
    greet("Dr. Esmaeili", "Greetings to")
end
\end{salamcode}

\begin{salamout}
Hello Sara
Greetings to Dr. Esmaeili
\end{salamout}

Parameters with default values must come after the ordinary parameters.

\section{The scope of variables}

A variable declared inside a function exists only inside that function. This
property is called \textbf{scope}. \code{main} cannot see the variables
inside \code{greet}, and the other way round. Parameters are the same: boxes
that are made when the function starts and disappear when it ends.

The more important point is that numbers, strings and boolean values reach a
function as \emph{copies}. Whatever the function does with its own copy, the
caller's variable stays untouched (arrays are the exception; see chapter
\ref{ch:arrays}):

\program{A function receives a copy}
\begin{salamcode}
func double(n: int): int:
    mut result := n
    result *= 2
    ret result
end

func main:
    value := 21
    println "Result of the function:", double(value)
    println "Original value:", value
end
\end{salamcode}

\begin{salamout}
Result of the function: 42
Original value: 21
\end{salamout}

This is the same thing we saw in chapter \ref{ch:variables} about values
being copied. If a function wants to give a result to the outside world, the
right way is \code{ret}, not tampering with its argument.

\subsection{Global constants}

Sometimes a value is needed in several functions. Constants declared with
\code{const} at the top of the program, outside every function, are
available everywhere:

\program{A global constant}
\begin{salamcode}
const tax := 9

func final price(price: int): int:
    ret price + ((price * tax) / 100)
end

func main:
    println final price(100)
    println final price(250)
    println "Rate:", tax, "percent"
end
\end{salamcode}

\begin{salamout}
109
272
Rate: 9 percent
\end{salamout}

Outside functions, declare only \code{const}, never a changeable variable. A
global variable that any function can change makes faults very hard to
track down; it is better for every function to receive what it needs through
parameters and to hand back what it produces through its return value.

\begin{deep}[A function that calls itself]
A function may call itself from inside its own body. This is called
\textbf{recursion}, and interesting things can be done with it. The program
below counts down using exactly this trick: each call sends a number one
smaller, until it reaches zero and the chain stops. A \code{ret} with no
value ends a function that has no result on the spot. We will not go further
into this subject in this book, but it is good to know that it exists.
\end{deep}

\program{A recursive countdown}
\begin{salamcode}
func count down(n: int):
    if n == 0:
        println "Done"
        ret
    end
    println n
    count down(n - 1)
end

func main:
    count down(3)
end
\end{salamcode}
\begin{salamout}
3
2
1
Done
\end{salamout}

\section{A more complete program}

A restaurant bill in which every part of the job is handed to a function:

\program{A restaurant bill}
\begin{salamcode}
const service := 10

func order amount(unit price: int, count: int): int:
    ret unit price * count
end

func apply discount(amount: int, percent: int): int:
    ret amount - ((amount * percent) / 100)
end

func show row(label: str, amount: int):
    println label + ":", amount, "dollars"
end

func main:
    kebab := order amount(24, 2)
    soda := order amount(3, 4)
    total := kebab + soda
    service charge := (total * service) / 100

    show row("Kebab", kebab)
    show row("Soda", soda)
    show row("Total", total)
    show row("Service charge", service charge)
    show row("Amount due", apply discount(total + service charge, 5))
end
\end{salamcode}

\begin{salamout}
Kebab: 48 dollars
Soda: 12 dollars
Total: 60 dollars
Service charge: 6 dollars
Amount due: 63 dollars
\end{salamout}

Read the \code{main} of this program: it is almost plain English, because
every detail of the arithmetic is hidden inside a function with a
descriptive name. Look at the last line: the result of
\code{apply discount} is passed straight in as the argument of
\code{show row}. This is how functions fit inside one another.

\begin{summary}
\begin{itemize}
  \item A function is a named block: \code{func name(parameters): return type:}
        \ldots\ \code{end}. It is called with its name and parentheses.
  \item Functions are defined at the top level of the program; their order
        does not matter.
  \item Parameters are written with their types: \code{name: str}. The
        arguments reach them in order; numbers and strings as copies.
  \item \code{ret value} hands the result back and finishes the function.
        A call to such a function is itself a value.
  \item A parameter may have a default value.
  \item Variables inside a function exist only in that function. For
        shared values, declare a global \code{const}.
\end{itemize}
\end{summary}

\section*{Exercises}
\markright{Exercises}

\exercise Write a function called \code{maximum} that takes two integers
and returns the larger one. Then use it to find the largest of three
numbers.

\exercise Write a function that takes the radius of a circle (floating-point)
and returns its area (take pi as 3.1416). With a loop, print the areas of
circles with radii 1 to 5.

\exercise Write a function called \code{is leap year} that takes a year and
says whether it is a leap year. The rule: a year is a leap year if it is
divisible by 4, except that a year divisible by 100 is a leap year only when
it is also divisible by 400. Then print the result for the years 2023 to
2028.

\exercise Write a function that takes a number of seconds and returns it as
a string of the form \code{"1:02:05"} (hours:minutes:seconds). Hint: use the
calculation from chapter \ref{ch:operators} and join strings together.

\exercise Rewrite the multiplication table program of chapter
\ref{ch:loops} with a function called \code{table row} that takes the row
number and the number of columns and prints one row of the table. Give the
number of columns a default value of 10.

\exercise The function below has a fault. Find it without running the
program:
\begin{salamcode}
func absolute(n: int): int:
    if n < 0:
        ret -n
    else n > 0:
        ret n
    end
end
\end{salamcode}
